% Source: https://artofproblemsolving.com/community/c3h2050770_help_needed_with_polynomial_problem
% Stable identity: aops-topic-2050770-q01
Suppose $P(p)=(p^3+2p+m)Q(p)$. Because the divisor is monic with integer coefficients, long division shows that $Q$ has integer coefficients. Comparing constant terms gives $m\cdot Q(0)=30$, so $m\mid30$. Substituting $p=1$ gives $(3+m)\,Q(1)=P(1)=66$, so $m+3$ divides $66$. The divisors of $66$ that are greater than $3$ are $6,11,22,33,66$, so $m\in\{3,8,19,30,63\}$, and with $m\mid30$ we get $m\in\{3,30\}$. Substituting $p=2$ gives $(12+m)\mid P(2)=5280$. For $m=30$, $42\nmid5280$, so $m=30$ fails.

For $m=3$, the division is exact:
\[P(p)=(p^3+2p+3)\left(p^9-p^8+p^7-p^6+p^5-p^4+p^3-p^2+p+10\right),\]
which can be checked by expanding. Hence $m=\boxed{3}$.
